\answerAnchor{MATH1-CALC-0015}
\begin{problemMeta}
\textbf{题目编号：} \texttt{MATH1-CALC-0015}\quad
\textbf{对应题面：} \problemRef{MATH1-CALC-0015}
\end{problemMeta}

\begin{solutionBox}
\textbf{A　指数因子。}

目标式是 \(f'+2f=0\)，注意到
\[
\bigl(e^{2x}f(x)\bigr)'=e^{2x}\bigl(f'(x)+2f(x)\bigr).
\]
令 \(F(x)=e^{2x}f(x)\)。由题设，
\[
F\in C[0,1]\cap D(0,1),\qquad F(0)=F(1)=0.
\]
由罗尔定理，存在 \(\xi\in(0,1)\) 使
\[
0=F'(\xi)=e^{2\xi}\bigl(f'(\xi)+2f(\xi)\bigr).
\]
因为 \(e^{2\xi}>0\)，所以 \(f'(\xi)+2f(\xi)=0\)。

\textbf{核验：}取 \(f(x)=x(1-x)e^{-2x}\)，题设成立；目标式等价于
\((x(1-x))'=0\)，可取 \(\xi=1/2\)。
\end{solutionBox}

\begin{solutionBox}
\textbf{B　除以幂函数。}

目标式移项为 \(xf'-2f=0\)，而
\[
\left(\frac{f(x)}{x^2}\right)'
=\frac{xf'(x)-2f(x)}{x^3}.
\]
由于区间为 \([1,2]\)，分母不为零。令
\[
F(x)=\frac{f(x)}{x^2}.
\]
则 \(F\in C[1,2]\cap D(1,2)\)，并且
\[
F(1)=f(1),\qquad F(2)=\frac{f(2)}4=f(1).
\]
由罗尔定理，存在 \(\xi\in(1,2)\) 使 \(F'(\xi)=0\)。因
\(\xi^3>0\)，故
\[
\xi f'(\xi)-2f(\xi)=0.
\]

\textbf{核验：}取 \(f(x)=Cx^2\)，端点条件成立，目标式在每一点成立。
\end{solutionBox}

\begin{solutionBox}
\textbf{C　带权积分构造。}

令
\[
F(x)=\int_0^x 2t\bigl(f(t)-I\bigr)\,\mathrm dt.
\]
因 \(f\) 连续，\(F\in C[0,1]\cap D(0,1)\)，且
\[
F(0)=0,
\]
\[
\begin{aligned}
F(1)
&=\int_0^1 2t f(t)\,\mathrm dt
  -I\int_0^1 2t\,\mathrm dt\\
&=I-I=0.
\end{aligned}
\]
由罗尔定理，存在 \(\xi\in(0,1)\) 使
\[
0=F'(\xi)=2\xi\bigl(f(\xi)-I\bigr).
\]
这里 \(\xi\in(0,1)\)，所以 \(2\xi>0\)，进而 \(f(\xi)=I\)。

\textbf{核验：}权函数 \(2x\ge0\) 且积分为 \(1\)，故 \(I\) 是
\(f\) 的加权平均；结论与连续函数介值性一致。上面的构造还说明了
为何必须利用 \(\xi>0\) 排除乘子为零。
\end{solutionBox}

\begin{solutionBox}
\textbf{M　参数插值与两次罗尔定理。}

构造通过三个已知点的二次多项式
\[
p(x)=\frac{\lambda}{2}x(3-x).
\]
确有
\[
p(0)=p(3)=0,\qquad p(1)=\lambda,\qquad p''(x)=-\lambda.
\]
令 \(F(x)=f(x)-p(x)\)，则
\[
F(0)=F(1)=F(3)=0.
\]
分别在 \([0,1]\) 和 \([1,3]\) 上使用罗尔定理，存在
\[
u\in(0,1),\qquad v\in(1,3)
\]
使 \(F'(u)=F'(v)=0\)。由于 \(f\) 在 \((0,3)\) 内二阶可导，
\(F'\) 在 \([u,v]\) 上连续、在 \((u,v)\) 内可导。再用一次罗尔
定理，存在 \(\xi\in(u,v)\subset(0,3)\) 使
\[
0=F''(\xi)=f''(\xi)-p''(\xi)=f''(\xi)+\lambda.
\]
所以 \(f''(\xi)=-\lambda\)。

\textbf{退化检查：}当 \(\lambda=0\) 时，\(p\equiv0\)，结论退化为
三个零点推出某点 \(f''(\xi)=0\)，上述证明仍然成立。取 \(f=p\)
时，目标式处处成立，也核验了常数和符号。
\end{solutionBox}

\answerAnchor{MATH1-CALC-0016}
\begin{problemMeta}
\textbf{题目编号：} \texttt{MATH1-CALC-0016}\quad
\textbf{对应题面：} \problemRef{MATH1-CALC-0016}
\end{problemMeta}

\begin{solutionBox}
\textbf{A　先判型，再决定能否使用。}

\((1)\) 分子、分母均趋于 \(0\)，是 \(0/0\) 型，可直接使用：
\[
\lim_{x\to0}\frac{1-\cos x}{x^2}
=\lim_{x\to0}\frac{\sin x}{2x}
=\frac12.
\]

\((2)\) 是 \(0\cdot(-\infty)\) 型，不能直接使用。先改写为
\[
x\ln x=\frac{\ln x}{1/x},
\]
此时是 \((-\infty)/(+\infty)\) 型，故
\[
\lim_{x\to0^+}\frac{\ln x}{1/x}
=\lim_{x\to0^+}\frac{1/x}{-1/x^2}
=\lim_{x\to0^+}(-x)=0.
\]

\((3)\) 不是未定式。右极限为 \(+\infty\)，左极限为 \(-\infty\)，
所以两侧极限不存在，不能用洛必达法则。

\((4)\) 是 \(\infty/\infty\) 型，可直接使用：
\[
\lim_{x\to+\infty}\frac{\ln x}{x}
=\lim_{x\to+\infty}\frac{1/x}{1}=0.
\]
\end{solutionBox}

\begin{solutionBox}
\textbf{B　一次洛必达后做代数化简。}

原式是 \(0/0\) 型：
\[
\begin{aligned}
\lim_{x\to0}\frac{\ln(1+x)-x}{x^2}
&=\lim_{x\to0}\frac{\frac1{1+x}-1}{2x}\\
&=\lim_{x\to0}\frac{-x/(1+x)}{2x}
=-\frac12.
\end{aligned}
\]
\textbf{核验：}
\(\ln(1+x)=x-x^2/2+o(x^2)\)，故结果一致。
\end{solutionBox}

\begin{solutionBox}
\textbf{C　先化乘积，再用基本极限。}

\[
\begin{aligned}
\frac{\tan x-\sin x}{x^3}
&=\frac{\sin x(1-\cos x)}{x^3\cos x}\\
&=\frac{\sin x}{x}\cdot
  \frac{1-\cos x}{x^2}\cdot
  \frac1{\cos x}.
\end{aligned}
\]
三个因子分别趋于 \(1,1/2,1\)，所以
\[
\boxed{\lim_{x\to0}\frac{\tan x-\sin x}{x^3}=\frac12}.
\]
这里直接结构化简比连续使用洛必达法则更短。
\end{solutionBox}

\begin{solutionBox}
\textbf{M　原比值有极限，不迫使导数比有极限。}

\[
\frac{f(x)}{g(x)}
=1+x\sin\frac1{x^2}.
\]
由
\[
\left|x\sin\frac1{x^2}\right|\le |x|\to0,
\]
得 \(f(x)/g(x)\to1\)。

另一方面，
\[
f'(x)
=1+2x\sin\frac1{x^2}
-\frac2x\cos\frac1{x^2},
\qquad g'(x)=1.
\]
取
\[
x_n=\frac1{\sqrt{2\pi n}},
\]
则正弦项为 \(0\)、余弦项为 \(1\)，故
\[
\frac{f'(x_n)}{g'(x_n)}=1-\frac2{x_n}\to-\infty.
\]
再取
\[
y_n=\frac1{\sqrt{(2n+1)\pi}},
\]
则正弦项为 \(0\)、余弦项为 \(-1\)，故
\[
\frac{f'(y_n)}{g'(y_n)}=1+\frac2{y_n}\to+\infty.
\]
所以导数比极限不存在。这反驳了错误命题
\[
\lim\frac fg\text{ 存在}
\quad\Longrightarrow\quad
\lim\frac{f'}{g'}\text{ 存在且与其相等}.
\]
洛必达法则是充分条件工具，不具有上述逆命题。
\end{solutionBox}

\answerAnchor{MATH1-CALC-0017}
\begin{problemMeta}
\textbf{题目编号：} \texttt{MATH1-CALC-0017}\quad
\textbf{对应题面：} \problemRef{MATH1-CALC-0017}
\end{problemMeta}

\begin{solutionBox}
\textbf{A　把 \(0\cdot\infty\) 改写为 \(0/0\)。}

\[
x\left(e^{1/x}-1\right)
=\frac{e^{1/x}-1}{1/x}.
\]
此时分子、分母均趋于 \(0\)，故
\[
\lim_{x\to+\infty}
\frac{e^{1/x}-1}{1/x}
=
\lim_{x\to+\infty}
\frac{-e^{1/x}/x^2}{-1/x^2}
=1.
\]
\textbf{核验：}令 \(t=1/x\to0^+\)，原式化为
\((e^t-1)/t\to1\)。
\end{solutionBox}

\begin{solutionBox}
\textbf{B　根式差先乘共轭式。}

\[
\begin{aligned}
&\sqrt{x^2+x+1}-\sqrt{x^2-x+1}\\
&=\frac{2x}
{\sqrt{x^2+x+1}+\sqrt{x^2-x+1}}\\
&=\frac{2}
{\sqrt{1+1/x+1/x^2}+\sqrt{1-1/x+1/x^2}}.
\end{aligned}
\]
因为 \(x\to+\infty\)，提出根号中的 \(x^2\) 时
\(\sqrt{x^2}=x\)。令 \(x\to+\infty\)，得
\[
\boxed{\lim=1}.
\]
\end{solutionBox}

\begin{solutionBox}
\textbf{C　幂指型先取对数。}

令
\[
y=\left(\frac{x+1}{x-1}\right)^x.
\]
当 \(x>1\) 时底数为正，并且
\[
\ln y
=x\ln\frac{1+1/x}{1-1/x}.
\]
令 \(t=1/x\to0^+\)，则
\[
\lim\ln y
=\lim_{t\to0^+}\frac{\ln(1+t)-\ln(1-t)}{t}.
\]
由洛必达法则，
\[
\lim_{t\to0^+}
\left(\frac1{1+t}+\frac1{1-t}\right)=2.
\]
所以
\[
\boxed{\lim_{x\to+\infty}y=e^2}.
\]
\textbf{核验：}利用
\(\ln(1+t)-\ln(1-t)=2t+o(t)\)，同样得到对数极限为 \(2\)。
\end{solutionBox}

\begin{solutionBox}
\textbf{M　参数决定对数主部的阶数。}

令
\[
y_\alpha=\left(1+x^{-\alpha}\right)^x.
\]
底数恒正，且
\[
\ln y_\alpha
=x\ln\left(1+x^{-\alpha}\right)
=x^{1-\alpha}
\cdot\frac{\ln(1+x^{-\alpha})}{x^{-\alpha}}.
\]
其中第二因子趋于 \(1\)。于是
\[
\lim\ln y_\alpha=
\begin{cases}
+\infty,&0<\alpha<1,\\
1,&\alpha=1,\\
0,&\alpha>1.
\end{cases}
\]
指数还原后，
\[
\boxed{
L_\alpha=
\begin{cases}
+\infty,&0<\alpha<1,\\
e,&\alpha=1,\\
1,&\alpha>1.
\end{cases}}
\]
\end{solutionBox}

\answerAnchor{MATH1-CALC-0018}
\begin{problemMeta}
\textbf{题目编号：} \texttt{MATH1-CALC-0018}\quad
\textbf{对应题面：} \problemRef{MATH1-CALC-0018}
\end{problemMeta}

\begin{solutionBox}
\textbf{A1　由各阶导数确定系数。}

令 \(f(x)=(1+x)^{1/2}\)。则
\[
f(0)=1,\qquad
f'(0)=\frac12,
\]
并且
\[
f''(x)=-\frac14(1+x)^{-3/2},
\qquad f''(0)=-\frac14.
\]
所以
\[
\boxed{
\sqrt{1+x}
=1+\frac{x}{2}-\frac{x^2}{8}+o(x^2)
}.
\]
\end{solutionBox}

\begin{solutionBox}
\textbf{A2　二项式系数与复合小量。}

使用
\[
(1+u)^\alpha
=1+\alpha u
+\frac{\alpha(\alpha-1)}{2}u^2
+\frac{\alpha(\alpha-1)(\alpha-2)}{6}u^3
+o(u^3),
\]
取
\[
\alpha=-\frac12,\qquad u=-2x.
\]
逐项计算得
\[
\boxed{
(1-2x)^{-1/2}
=1+x+\frac32x^2+\frac52x^3+o(x^3)
}.
\]
\end{solutionBox}

\begin{solutionBox}
\textbf{B1　相消到五阶，必须保留五阶项。}

\[
\sin x
=x-\frac{x^3}{6}+\frac{x^5}{120}+o(x^5).
\]
所以
\[
\sin x-x+\frac{x^3}{6}
=\frac{x^5}{120}+o(x^5),
\]
从而
\[
\boxed{\lim=\frac1{120}}.
\]
\end{solutionBox}

\begin{solutionBox}
\textbf{B2　比较两个奇函数的三阶主部。}

\[
\arctan x=x-\frac{x^3}{3}+o(x^3),
\qquad
\sin x=x-\frac{x^3}{6}+o(x^3).
\]
相减得
\[
\arctan x-\sin x
=-\frac{x^3}{6}+o(x^3),
\]
故
\[
\boxed{\lim=-\frac16}.
\]
\end{solutionBox}

\begin{solutionBox}
\textbf{C　复合小量必须展开幂次中的交叉项。}

令
\[
u=x+x^2.
\]
则
\[
\ln(1+u)=u-\frac{u^2}{2}+\frac{u^3}{3}+o(u^3),
\]
且
\[
u^2=x^2+2x^3+o(x^3),
\qquad
u^3=x^3+o(x^3).
\]
所以
\[
\begin{aligned}
\ln(1+x+x^2)
&=x+x^2-\frac12(x^2+2x^3)+\frac{x^3}{3}+o(x^3)\\
&=x+\frac{x^2}{2}-\frac{2x^3}{3}+o(x^3).
\end{aligned}
\]
因此
\[
\boxed{\lim=-\frac23}.
\]
\end{solutionBox}

\begin{solutionBox}
\textbf{M　参数匹配与拉格朗日余项。}

\textbf{(1) 参数匹配。}
\[
\ln(1+x)
=x-\frac{x^2}{2}+\frac{x^3}{3}-\frac{x^4}{4}+o(x^4).
\]
所以分子为
\[
\left(a-\frac12\right)x^2
+\left(b+\frac13\right)x^3
-\frac{x^4}{4}+o(x^4).
\]
极限有限要求低于四阶的系数全部为零，即
\[
a=\frac12,\qquad b=-\frac13.
\]
此时
\[
\boxed{\lim=-\frac14}.
\]

\textbf{(2) 余项控符号与大小。}

当 \(x=0\) 时两端均取等号，结论显然。下设
\(0<x\le\pi/2\)。把 \(\cos x\) 在 \(0\) 处展开到三阶，
则存在 \(\xi\in(0,x)\)，使
\[
\cos x
=1-\frac{x^2}{2}+\frac{\cos\xi}{4!}x^4.
\]
此时
\[
0\le\xi\le\frac{\pi}{2},
\qquad 0\le\cos\xi\le1.
\]
于是
\[
0\le\frac{\cos\xi}{24}x^4\le\frac{x^4}{24},
\]
从而
\[
\boxed{
1-\frac{x^2}{2}
\le\cos x
\le1-\frac{x^2}{2}+\frac{x^4}{24}
}.
\]
\end{solutionBox}
